%%%PAGE 83 rot=0 legibility=fair summary=函数分析(物理量的变化率与面积)；定义式/决定式；电势能系统；磁场中I与粒子不受力条件、电动势E=W/q；功能关系汇总；能量图象(无f/有f、斜面上滑下滑)；ε=Bl√(2ax)匀加速ε-x图；机车启动(恒P/恒a)
%%%BLOCK id=083-01 ch=99 sec=页边零星·化学 kind=misc alt=none cont=none y=0.03-0.08
% 页面上缘右侧页眉“No.”旁有零星字迹(与p081上缘相同，疑为化学·气体密度)，照录
$\rho=1.429\unit{g\cdot L^{-1}}$
%%%END
%%%BLOCK id=083-02 ch=09 sec=电势能·系统 kind=formula alt=06 cont=none y=0.04-0.20
% 红笔：三个q，边上标l的小斜线(三角形三顶点)
%FIG p083_a 0.33 0.04 0.465 0.1155
\notefig[0.25]{figures/p083_a.png}
\red{$E_p=\dfrac12\,\dfrac{3kq}{r}\,q$}

$q$与$Q$系统电势能$=E_{pq}$
%%%END
%%%BLOCK id=083-03 ch=11 sec=磁感应强度·定义式 kind=formula alt=none cont=none y=0.07-0.12
\[
B=\frac{F}{IL\sin\theta}
\]
%%%END
%%%BLOCK id=083-04 ch=05 sec=重力场·场强 kind=note alt=09 cont=none y=0.10-0.17
% 原稿：g的箭头向下指向小圆圈“○”
$\downarrow g$\quad $\circ$\quad 重力场\qquad $g=\dfrac{G}{m}$

线与负电荷一样
%%%END
%%%BLOCK id=083-05 ch=09 sec=电通量 kind=formula alt=none cont=none y=0.07-0.11
\[
\Phi_{\text{电}}=E_{\perp}S
\]
%%%END
%%%BLOCK id=083-06 ch=06 sec=机械能·能量关系 kind=formula alt=none cont=none y=0.05-0.12
$Q=E_2-E_1$\qquad $E_k=\dfrac12mv^2$

$E_{\text{机}}=E_k+E_p$
%%%END
%%%BLOCK id=083-07 ch=90 sec=定义式与决定式 kind=summary alt=01,02,03,09,10,11 cont=none y=0.11-0.23
% 原稿左侧有大括号：{定义式 比值 (无关)；决定式 (相关)}，“无关”“相关”均被圈起；公式按原稿行序排列
\[
\left\{\begin{aligned}
&\text{定义式\ \ 比值}\\[8pt]
&\text{决定式}
\end{aligned}\right.
\]

$v=\dfrac{\Delta x}{\Delta t}$\quad $k=\dfrac{F}{x}$\quad $E=\dfrac{F}{q}$\quad $\unclear{R=\dfrac{U}{I}}$ % CHECK: 末式笔迹形似“k=v/L”，按R=U/I转录(本页“R”常写得像“k”)

$I=\dfrac{q}{t}$\quad $C=\dfrac{Q}{U}$\quad $\Phi=B_{\perp}S$

\fbox{无关}\quad $p=mv$\quad $I=Ft$

\fbox{相关}\quad $a=\dfrac{F}{m}$\quad $E=\dfrac{kQ}{r^2}$

$I=nesv$\quad $R=\dfrac{\rho L}{S}$\quad $C=\dfrac{\varepsilon S}{4\pi kd}$
%%%END
%%%BLOCK id=083-08 ch=90 sec=函数分析·变化率与面积 kind=formula alt=01,06,09,10,12 cont=none y=0.12-0.40
\textbf{函数分析}

% 原稿左侧红笔大括号括住下列各式
\[
\left\{\begin{aligned}
a&=\frac{\Delta v}{\Delta t}\\[6pt]
E_{\text{\unclear{感}}}&=-n\frac{\Delta\Phi}{\Delta t}\qquad E_{\text{场}}=-\red{\frac{\Delta\varphi}{\Delta x}}\\[6pt]
i&=\frac{\Delta q}{\Delta t}\\[6pt]
F&=\frac{\Delta p}{\Delta t}\\[6pt]
P&=\frac{\Delta W}{\Delta t}\\[6pt]
F_{\text{合}}&=\frac{\Delta E_k}{\Delta x}\\[6pt]
F_{\text{其他}}&=\frac{\Delta E_{\text{机}}}{\Delta x}\\[6pt]
R&=\frac{U_1}{I_1}=\frac{U_2}{I_2}\neq\frac{U_2-U_1}{I_2-I_1}\\[6pt]
P&=U_1I_1\qquad\text{面积}\ S=\int xy
\end{aligned}\right.
\]
% 原稿E_场式中黑笔旧写的分式被红笔涂划，右侧红笔重写为Δφ/Δx；E_感下标笔迹似“源/感”

$\dfrac{\Delta v}{\Delta t}$\quad 匀圆变化\quad 方向变
%%%END
%%%BLOCK id=083-09 ch=06 sec=重力势能·系统 kind=knowledge alt=none cont=none y=0.20-0.23
$m$与地球系统重力势能$=E_{pm}$
%%%END
%%%BLOCK id=083-10 ch=11 sec=磁场·受力条件与方向 kind=note alt=none cont=none y=0.24-0.37
% 原稿此块及下块左侧有红笔大括号
在$B$中放$I$\quad 测不了$B$\quad $B=\dfrac{F}{IL\sin\theta}$\quad $\theta$未知.

$I$不受磁场力\quad ①$B=0$\quad ②$I\parallel B$\qquad 粒子：①$B=0$\quad ②$v\parallel B$

电流受磁场力方向与磁场垂直\qquad 三垂直\quad $F=qv\times B$\quad $F=IL\times B$

\unclear{场}真\quad 线假
%%%END
%%%BLOCK id=083-11 ch=10 sec=电动势·非静电力 kind=knowledge alt=none cont=none y=0.37-0.40
% 以下三处原稿有红笔下划线
$E=\dfrac{W}{q}$\quad \underline{$W$为非静电力做的功}\quad \underline{$E$非静电力做功本领}\quad \underline{不是蓄电}
%%%END
%%%BLOCK id=083-12 ch=06 sec=功能关系 kind=formula alt=05,09,12 cont=none y=0.41-0.60
\textbf{功能关系}

% 原稿左侧红笔大括号；“=−mgΔh”“=−kQq/r”写在W_G、W_电式的上方右侧
\[
\left\{\begin{aligned}
&W_{\text{合}}=\Delta E_k\qquad W_G=-\Delta E_{pG}=-mg\Delta h\qquad W_{\text{电}}=-\Delta E_{p\text{电}}=-\frac{kQq}{r}\\[6pt]
&W_{\text{弹}}=-\Delta E_{p\text{弹}}=-\frac12k\Delta x^2\\[6pt]
&W_{\text{引}}=-\Delta E_{p\text{引}}=\frac{GMm}{r} % CHECK: 下标笔迹似“31”，按“引”(万有引力)转录
\\[6pt]
&W_f=-Q_f\\[6pt]
&W_{\text{安}}=-Q_{\text{电}}\\[6pt]
&W_{\text{其他}}=\Delta E_{\text{机}}
\end{aligned}\right.
\]

\[
W_F+W_G+W_{\text{电}}+W_{\text{弹}}+W_f=\Delta E_k
\]
%%%END
%%%BLOCK id=083-13 ch=06 sec=能量图象·有无阻力 kind=method alt=01 cont=none y=0.40-0.58
% 左上一小图：“O”(小球)，下方“↓ ↑”箭头与横线
%FIG p083_b 0.055 0.42 0.16 0.52
\notefig[0.25]{figures/p083_b.png}
%FIG p083_c 0.195 0.405 0.61 0.552
\notefig[0.55]{figures/p083_c.png}
无$f$时

$E_{k1}=mgh$\qquad\qquad 有$f$时\quad 斜率变化
%%%END
%%%BLOCK id=083-14 ch=06 sec=能量图象·斜面上滑下滑 kind=method alt=03 cont=none y=0.58-0.71
%FIG p083_d 0.085 0.589 0.875 0.6805
\notefig[0.85]{figures/p083_d.png}
先\unclear{加}后减速\unclear{来}\qquad $E_k$开口全向上，$E_p$开口全向下。$E_{\text{机}}$-$X$不分段 % CHECK: 首句字迹潦草，“先加(口)后减速来”照录
%%%END
%%%BLOCK id=083-15 ch=12 sec=电磁感应·图象 kind=method alt=01 cont=none y=0.71-0.81
% 原稿“□→”为一个小方块加向右箭头
$\square\to$\quad $\varepsilon=Bl\sqrt{2ax}$\quad \red{匀加速}
%FIG p083_f 0.335 0.71 0.52 0.81
\notefig[0.3]{figures/p083_f.png}
%%%END
%%%BLOCK id=083-16 ch=06 sec=机车启动 kind=method alt=03 cont=none y=0.85-1.00
\textbf{机车启动}

① 恒$P$启动\qquad $\dfrac{P}{v}-f=ma$\qquad $V_m=\dfrac{P}{f}$
%FIG p083_g 0.125 0.962 0.27 1.0
\notefig[0.25]{figures/p083_g.png}
%FIG p083_h 0.37 0.898 0.525 0.99
\notefig[0.3]{figures/p083_h.png}

② 恒$a$启动
\[
\left\{\begin{aligned}
&F-f=ma\\
&P=Fv
\end{aligned}\right.
\]
\[
P=ma+f\cdot\unclear{at}
\]
% CHECK: 末行被页面下缘截断，仅见“P=ma+f·?t”，原稿似无括号
%FIG p083_i 0.745 0.88 0.935 0.995
\notefig[0.3]{figures/p083_i.png}
%%%END
%%%PAGEEND 83
